% ============================ RELATED WORK ============================
% Structure: manual + three autonomy families (reactive/geometry-based, hierarchical, RL).
Flipper control ranges from manual control to three families of autonomy: reactive and
geometry-based control, hierarchical control, and reinforcement learning.
Table~\ref{tab:sensing} summarizes the compared controllers, their sensing requirements,
and their traversal quality.

\subsection{Manual control}
In direct teleoperation the operator drives the flippers by hand~\cite{cihala2025}, or
switches between preset flipper modes~\cite{larochelle2013,azayev2022}. Manual control
needs no map: the operator observes the terrain, closes the control loop by hand, and also stabilizes
the robot.

\subsection{Reactive and geometry-based control}
A family that computes flipper commands from an explicit terrain model, reactively~\cite{cihala2025}
or by planning, and is the \emph{most} map-dependent of all. Chen et al. combine dynamic programming with a
fast geometry-based pose predictor over the elevation
point cloud~\cite{chen2023}; Yuan et al. plan in the flipper configuration space
from known stair geometry with motion-capture ground truth~\cite{yuan2019}; Oehler
and von Stryk predict pose on a signed-distance field built from
LiDAR~\cite{oehler2024}; and Xu et al. solve a receding-horizon hybrid trajectory
optimization over a simplified terrain sampled from an a priori 3D occupancy map with
FAST-LIO2 odometry~\cite{xu2024}; and, most recently, Zhang et al. plan the path and flipper
maneuvers with RRT* over a contact-model pose predictor, using learned semantic
segmentation of the elevation map only as a perception front-end~\cite{zhang2025}.

\subsection{Hierarchical control}
A family that composes low-level behaviors under a selector. The adaptive-traversability line~\cite{zimmermann2014} classifies
terrain into discrete flipper modes from an elevation map, and for the selected mode a
simple low-level controller drives the flippers. A learned state machine~\cite{azayev2022}
switches among seven hand-crafted flipper morphologies selected from an elevation map. Each
state overlays a \emph{roll}-stabilization term and a stagnation-triggered
\emph{anti-stuck} term that lowers all four flippers to free the robot. It is the one prior
method that couples an anti-stuck term with attitude stabilization, and the origin of the
anti-stuck term we adopt. \TODO{keep this in related work?}

\subsection{Reinforcement learning}
FTR-Bench~\cite{ftrbench2025} standardizes this setting, benchmarking SAC, PPO, TD3, DDPG,
and TRPO on a flipper-track robot in Isaac Lab. Pecka et al. learn a
safety-constrained policy (Constrained REPS) whose state is body pitch plus the terrain
height read from an online octomap~\cite{pecka2016}. Discrete-action variants dominate the
flipper-specific literature: Pan et al. learn a dueling double-DQN over discrete
front/rear flipper increments~\cite{pan2023atd3qn} and extend it with an
intrinsic-curiosity module~\cite{pan2023icm}; both take a local elevation map,
downsampled from a LiDAR point-cloud map built with A-LOAM~\cite{loam2014,aloam}, plus
chassis pitch. Their reward favors smooth pitch only, and getting stuck is a terminal
penalty. Mitriakov et al.
learn staircase negotiation with a policy-gradient method, first from step-edge
geometry with marker-based ground truth~\cite{mitriakov2021} and then in an
open-source Gazebo/ROS framework from depth-image features with ground-truth
pose~\cite{mitriakov2022}, injecting stability through center-of-gravity and
pitch-rate reward shaping. Most recently, C-TRAC learns a contact-aware latent from noisy
LiDAR and drives an asymmetric-SAC policy that, unusually, stabilizes both roll and pitch
via a normalized-energy-stability-margin reward~\cite{ctrac2025}.

\subsection{Attitude stabilization and mapless posture control}\label{sec:related-posture}
Flipper-based attitude stabilization has been studied before, but nearly always on top of
an elevation map: a roll-only stabilization term inside a map-based state
machine~\cite{azayev2022}, and C-TRAC stabilizing both roll and pitch from a reward computed
on the map-based simulator state~\cite{ctrac2025}. Rocha et al.~\cite{rocha2023} regulate
roll, pitch, and ground clearance of an articulated tracked vehicle from the IMU and the
flipper joints alone, but the method is not designed for obstacle traversal. It needs no
exteroceptive sensor and no map, and it is the closest prior work to ours in what it
senses.

\subsection{Evaluating traversal quality}
We evaluate a controller by the quality of the traversal it produces, measured
on the robot's body: IMU-based smoothness such as pitch swing and vertical
shock~\cite{chen2023}, body attitude, and chassis-to-ground clearance.
Following~\cite{cihala2025,zimmermann2014} we score shock, tilt, and clearance
(Sec.~\ref{sec:metrics-block}).

\begin{table*}[t]
\centering
\caption{\textbf{Compared controllers and their traversal quality in Gazebo.} \emph{Map \&
loc.}: needs an onboard map. \emph{Stab.}: attitude axes with an
explicit stabilization objective (r roll, p pitch). \emph{Anti-stuck}: frees itself when
stuck. TQ: mean $\pm$ s.d.\ over ten runs of the twelve-obstacle course under the
ground-truth (GT) and the onboard map; $\Delta$: TQ change under the onboard map;
\emph{Override}: share of obstacle distance the \SAS{} override held the flippers. Manual
rides are driven by an operator at free speed. Best in bold, second best underlined.}
\label{tab:sensing}
\footnotesize
\setlength{\tabcolsep}{2.2pt}
\begin{tabular}{@{}llcccccccc@{}}
\toprule
Controller & Approach & Code & Map \& loc. & Stab. & Anti-stuck & \shortstack{TQ $\uparrow$\\GT map} & \shortstack{TQ $\uparrow$\\onboard map} & $\Delta$ $\uparrow$ & \shortstack{Override [\%] $\downarrow$\\GT / onboard} \\
\midrule
\multicolumn{10}{@{}l}{\emph{Manual}} \\
Manual control~\cite{cihala2025} & direct teleoperation & -- & \xmark & -- & \xmark & -- & 0.612 $\pm$ 0.100 & -- & -- \\
Manual modes~\cite{azayev2022} & operator switches preset modes & -- & \xmark & -- & \xmark & -- & 0.568 $\pm$ 0.101 & -- & -- \\
\midrule
\multicolumn{10}{@{}l}{\emph{Reactive and geometry-based}} \\
{\v C}\'ihala \GC/\OA~\cite{cihala2025} & clearance + obstacle-aware control & released & \cmark & -- & \xmark & \textbf{0.529} $\pm$ 0.011 & 0.486 $\pm$ 0.015 & $-0.043$ & 0.4 / 0.0 \\
Chen GFMP~\cite{chen2023} & pose prediction + dynamic prog. & reimpl. & \cmark & p & \xmark & 0.509 $\pm$ 0.010 & 0.481 $\pm$ 0.012 & $-0.027$ & 7.5 / 8.7 \\
Yuan~\cite{yuan2019} & configuration-space planning & reimpl. & \cmark & -- & \xmark & \underline{0.522} $\pm$ 0.009 & \textbf{0.498} $\pm$ 0.010 & $\underline{-0.024}$ & 11.0 / 8.2 \\
Oehler~\cite{oehler2024} & SDF pose prediction + margin search & released & \cmark & -- & \xmark & 0.499 $\pm$ 0.013 & 0.467 $\pm$ 0.013 & $-0.032$ & 9.9 / 7.1 \\
Xu HTO~\cite{xu2024} & hybrid trajectory optimization & reimpl. & \cmark & p & \xmark & 0.465 $\pm$ 0.007 & 0.439 $\pm$ 0.009 & $-0.026$ & 8.9 / 9.8 \\
Zhang~\cite{zhang2025} & contact-model pose prediction & reimpl. & \cmark & -- & \xmark & 0.493 $\pm$ 0.010 & 0.462 $\pm$ 0.012 & $-0.031$ & 12.7 / 15.2 \\
\midrule
\multicolumn{10}{@{}l}{\emph{Hierarchical}} \\
Azayev~\cite{azayev2022} & learned template state machine & released & \cmark & r & \cmark & 0.436 $\pm$ 0.008 & 0.406 $\pm$ 0.017 & $-0.030$ & 7.1 / 6.7 \\
\midrule
\multicolumn{10}{@{}l}{\emph{Reinforcement learning}} \\
Pecka~\cite{pecka2016} & safety-constrained REPS & reimpl. & \cmark & p & \xmark & 0.458 $\pm$ 0.010 & 0.434 $\pm$ 0.010 & $\underline{-0.024}$ & 7.8 / 9.7 \\
AT-D3QN~\cite{pan2023atd3qn} & dueling double-DQN & reimpl. & \cmark & p & \xmark & 0.473 $\pm$ 0.012 & 0.439 $\pm$ 0.012 & $-0.034$ & 11.4 / 12.7 \\
ICM-D3QN~\cite{pan2023icm} & D3QN + intrinsic curiosity & reimpl. & \cmark & p & \xmark & 0.466 $\pm$ 0.011 & 0.443 $\pm$ 0.009 & $\mathbf{-0.023}$ & 11.2 / 12.9 \\
Mitriakov~\cite{mitriakov2021} & policy-gradient stair climbing & released & \cmark & p & \xmark & 0.451 $\pm$ 0.009 & 0.408 $\pm$ 0.007 & $-0.044$ & 7.3 / 8.3 \\
C-TRAC~\cite{ctrac2025} & contact-aware asymmetric SAC & reimpl. & \cmark & r\&p & \xmark & 0.480 $\pm$ 0.013 & 0.398 $\pm$ 0.008 & $-0.082$ & 0.0 / 7.0 \\
PPO~\cite{sarga2026thesis} & PPO over a heightmap & ours & \cmark & r\&p & \xmark & 0.461 $\pm$ 0.008 & 0.438 $\pm$ 0.010 & $\mathbf{-0.023}$ & 11.7 / 13.7 \\
\midrule
\textbf{\SAS{} (ours)} & attitude stabilization + anti-stuck & ours & \xmark & r\&p & \cmark & \textbf{0.529} $\pm$ 0.008 & \underline{0.497} $\pm$ 0.013 & $-0.032$ & -- \\
\bottomrule
\end{tabular}
\end{table*}

